\documentclass[openany]{book}
\input{notes_white}
\usepackage{mathtools}
\settitle{Analysis Pre-class Notes}
\begin{document}
\title{\Large{\underemphasise{Algorithm Practice Problems}} \\ \vspace{5pt}
\large{A collection of problems to practice algorithms}}
\author{Eklavya Raman \\ \vspace{5pt} er24ms024@iiserkol.ac.in \\ \vspace{5pt}}
\date{\today}
\maketitle

\section{Exercises}
\subsection{Easy Problems}
\begin{enumerate}
\item \textbf{Majority Element} - Given an array of size $n$, find the element which appears the most number of times. If no such unique element exists, return -1.
\item \textbf{Two Sum} - Given an array of integers and a target sum, find two distinct indices such that the numbers at those indices add up to the target sum.
\item \textbf{Reverse a list} - Write a function that takes a list as input and returns the string reversed.
\end{enumerate}
\subsection{Intermediate problems}
\begin{enumerate}
    \item \textbf{Fibonacci Generation} - Write a liner and recursive program to generate the $n^{th}$ fibonacci number. 
    \item \textbf{List sort} - Write an algorithm which takes a list as an input and outputs the sorted array. 
    \item \textbf{Closest Pair} - Given a set of points in a 2D plane, find the closest pair of points.
    \item \textbf{Longest Common Subsequence} - Given two arrays, find the length of the longest subsequence present in both arrays.
\end{enumerate}
\subsection{Hard Problems}
These problems are harder and may require more advanced techniques or brute force methods. They are not expected to be solved in a short time, both come in the class NP, with knapsack being NP-complete.
\begin{enumerate}
    \item \textbf{Knapsack Problem} - Given a set of items, each with a weight and a value, determine the maximum value that can be obtained by selecting items without exceeding a given weight limit.
    \item \textbf{Sudoku Solver} - Write a program that solves a given Sudoku puzzle (puzzle given as a 3D array)
\end{enumerate}

\section{Solutions}
\begin{lstlisting}[language=Python]
# Majority Element
def majority_element(arr):
    count = {}
    for num in arr:
        if num in count:
            count[num] += 1
        else:
            count[num] = 1
    max_count = 0
    majority_element = -1
    for key, value in count.items():
        if value > max_count:
            max_count = value
            majority_element = key
        if value == max_count:
            majority_element = -1
    return majority_element

# Two Sum
def two_sum(arr, target):
    seen = {}
    for i, num in enumerate(arr):
        complement = target - num
        if complement in seen:
            return (seen[complement], i)
        seen[num] = i
    return None

# Reverse a list
def reverse_list(lst):
    return lst[::-1]

# Fibonacci Generation
# Linear Generator 
def fibonacci_linear(n):
    a, b = 0, 1
    for i in range(n):
        a, b = b, a + b
    return a
# Recursive Generator
def fibonacci_recursive(n):
    if n <= 1:
        return n
    return fibonacci_recursive(n - 1) + fibonacci_recursive(n - 2)

# List sort
def list_sort(lst):
    n = len(lst)
    for i in range(n):
        for j in range(0, n-i-1):
            if lst[j] > lst[j+1]:
                lst[j], lst[j+1] = lst[j+1], lst[j]
    return lst

# Closest Pair
import math
def closest_pair(points):
    def distance(p1, p2):
        return math.sqrt((p1[0] - p2[0]) ** 2 + (p1[1] - p2[1]) ** 2)

    min_dist = float('inf')
    pair = None
    for i in range(len(points)):
        for j in range(i + 1, len(points)):
            dist = distance(points[i], points[j])
            if dist < min_dist:
                min_dist = dist
                pair = (points[i], points[j])
    return pair

# Longest Common Subsequence
def longest_common_subsequence(arr1, arr2):
    n, m = len(arr1), len(arr2)
    max_len = max(n, m)
    longest = []
    for i in range(n):
        subseq = []
        for j in range(m):
            if arr1[i] == arr2[j]:
                subseq.append(arr1[i])
                for k in range(1, max_len - max(i, j)):
                    if arr1[i+k] == arr2[j+k]:
                        subseq.append(arr1[i+k])
                    else:
                        break
        if len(subseq) > len(longest):
            longest = subseq
    return longest


# Knapsack Problem
def knapsack(weights, values, capacity):
    def is_valid_combo(weights=weights, capacity=capacity):
        curr_weight = 0
        for i in weights:
            if curr_weight + i > capacity:
                return False
        return True
    import itertools
    max_val = 0
    for r in range(1, len(weights) + 1):
        for combo in itertools.combinations(range(len(weights)), r):
            combo_weights = [weights[i] for i in combo]
            combo_values = [values[i] for i in combo]
            if is_valid_combo(combo_weights):
                total_value = sum(combo_values)
                if total_value > max_val:
                    max_val = total_value
    return max_val


# Sudoku Solver
def sudoku_solver(board):
    def is_valid(board, row, col, num):
        for i in range(9):
            if board[row][i] == num or board[i][col] == num:
                return False
        start_row, start_col = 3 * (row // 3), 3 * (col // 3)
        for i in range(start_row, start_row + 3):
            for j in range(start_col, start_col + 3):
                if board[i][j] == num:
                    return False
        return True

    def solve(board):
        for row in range(9):
            for col in range(9):
                if board[row][col] == 0:
                    for num in range(1, 10):
                        if is_valid(board, row, col, num):
                            board[row][col] = num
                            if solve(board):
                                return True
                            board[row][col] = 0
                    return False
        return True

    solve(board)
    return board

\end{lstlisting}
\end{document}

